and the KZB local system has finite monodromy.
The ordered chiral chain complex of $V_k(\mathfrak{sl}_2)$
on a genus-$g$ curve $\Sigma_g$ therefore computes
a finite-dimensional invariant at each degree, and the
symmetric coinvariants recover the space of conformal
blocks (Tsuchiya--Ueno--Yamada~\cite{TUY89}).
The following proposition shows that the resulting dimension
is given by the Verlinde formula~\cite{Verlinde88}.

\begin{proposition}[Verlinde formula from ordered chiral
homology]
\label{prop:verlinde-from-ordered}
Let $k \geq 1$ be a positive integer, and let
$S_{jl} = \sqrt{2/(k+2)}\,\sin\!\bigl(\pi(j+1)(l+1)/(k+2)\bigr)$
be the modular $S$-matrix for $\widehat{\mathfrak{sl}}_2$ at
level~$k$, where $j, l \in \{0, 1, \ldots, k\}$.
\begin{enumerate}[label=\textup{(\roman*)}]
\item \textup{(Genus~$0$.)}
  The ordered chiral homology at genus~$0$ with no marked
  points recovers the unique vacuum:
  \begin{equation}\label{eq:verlinde-g0}
    \dim H^0\!\bigl(\cM_{0,0},\, V_k(\mathfrak{sl}_2)\bigr)
    = \sum_{j=0}^{k} S_{0j}^2 = 1.
  \end{equation}
  This is the unitarity of the $S$-matrix:
  $\sum_l S_{jl}^2 = \delta_{j0}$ for $j = 0$ follows from
  the orthogonality of the sine functions.

\item \textup{(Genus~$1$.)}
  The ordered chiral homology at genus~$1$ recovers the
  number of integrable representations:
  \begin{equation}\label{eq:verlinde-g1}
    \dim H^0\!\bigl(\cM_{1,0},\, V_k(\mathfrak{sl}_2)\bigr)
    = \sum_{j=0}^{k} S_{0j}^0 = k + 1.
  \end{equation}
  By the Zhu algebra identification $A(V_k(\mathfrak{sl}_2))
  \cong U(\mathfrak{sl}_2)/(e^{k+1}, f^{k+1})$, the $k+1$
  simple modules of $A(V_k)$ are $V_0, \ldots, V_k$, and
  the genus-$1$ conformal blocks are their characters.
  From the ordered chiral homology:
  Proposition~\textup{\ref{prop:ell-degree0}} computes
  the degree-$0$ center at integrable level as $\CC^{k+1}$
  \textup{(}the center of the integrable quotient of $Y_\hbar$
  at the root of unity $q = e^{2\pi i/(k+2)}$\textup{)}.

\item \textup{(General genus.)}
  At genus~$g \geq 2$, the Verlinde formula gives:
  \begin{equation}\label{eq:verlinde-general}
    Z_g(k) := \dim H^0\!\bigl(\cM_{g,0},\,
    V_k(\mathfrak{sl}_2)\bigr)
    = \sum_{j=0}^{k} S_{0j}^{2-2g}.
  \end{equation}
  This is an integer for all $g \geq 0$, $k \geq 1$.
  The first values are:
  \begin{center}
  \renewcommand{\arraystretch}{1.3}
  \begin{tabular}{@{} c r r r r r @{}}
  \toprule
  $k$ & $Z_0$ & $Z_1$ & $Z_2$ & $Z_3$ & $Z_4$ \\
  \midrule
  $1$ & $1$ & $2$ & $4$ & $8$ & $16$ \\
  $2$ & $1$ & $3$ & $10$ & $36$ & $136$ \\
  $3$ & $1$ & $4$ & $20$ & $120$ & $800$ \\
  $4$ & $1$ & $5$ & $35$ & $329$ & $3{,}611$ \\
  $5$ & $1$ & $6$ & $56$ & $784$ & $13{,}328$ \\
  \bottomrule
  \end{tabular}
  \end{center}
  At $k = 1$: $S_{00} = S_{01} = 1/\sqrt{2}$, so
  $Z_g(1) = 2 \cdot (1/\sqrt{2})^{2-2g} = 2^g$.

\item \textup{(Handle attachment.)}
  The ordered bar complex factorization under non-separating
  sewing recovers the TQFT handle-attachment formula:
  each isospin channel $j$ contributes a handle operator
  $H_j = S_{0j}^{-2}$, so
  \begin{equation}\label{eq:handle-attachment}
    Z_{g+1}
    = \sum_{j=0}^{k} S_{0j}^{-2} \cdot S_{0j}^{2-2g}
    = \sum_{j=0}^{k} S_{0j}^{2-2(g+1)}.
  \end{equation}

\item \textup{(Separating factorization.)}
  Under separating degeneration $\Sigma_g \rightsquigarrow
  \Sigma_{g_1} \cup \Sigma_{g_2}$ with $g = g_1 + g_2$,
  the bar complex factorization gives
  \begin{equation}\label{eq:sep-factorization}
    Z_g
    = \sum_{j=0}^{k} S_{0j}^{2-2g_1} \cdot S_{0j}^{2-2g_2}
    \cdot S_{0j}^{-2}
    = \sum_{j=0}^{k} S_{0j}^{2-2g}.
  \end{equation}
\end{enumerate}
\end{proposition}

\begin{proof}
\textup{(i)} The identity $\sum_j S_{0j}^2 = 1$ is the
$(0,0)$-entry of $S \cdot S^T = I$, which follows from the
orthogonality relation
$\sum_{j=0}^{k} \sin(\pi(j{+}1)a/(k{+}2))
\sin(\pi(j{+}1)b/(k{+}2)) = (k{+}2)\delta_{ab}/2$.
From the ordered chiral homology: at genus~$0$ with no
marked points, the conformal block space is the space of
vacua $\CC \cdot |0\rangle$, which is one-dimensional.

\textup{(ii)} At genus~$1$, $S_{0j}^0 = 1$ for all $j$,
so $Z_1 = k + 1$. This has three independent derivations:
\begin{itemize}
\item \textup{(S-matrix.)} Direct computation:
  $\sum_{j=0}^{k} 1 = k + 1$.
\item \textup{(Zhu algebra.)} The Zhu algebra
  $A(V_k(\mathfrak{sl}_2))$ has $k+1$ simple modules;
  each contributes a one-dimensional space of genus-$1$
  conformal blocks (the character).
\item \textup{(Ordered center.)}
  Proposition~\ref{prop:ell-degree0} identifies the
  degree-$0$ center of $Y_\hbar(\mathfrak{sl}_2)$ at
  integrable level $k$ as $\CC^{k+1}$ (the center of
  the integrable quotient). The averaging map $\av_0$ is
  an isomorphism (nothing to symmetrise at degree~$0$),
  so all $k+1$ dimensions survive.
\end{itemize}

\textup{(iii)} The Verlinde formula~\cite{Verlinde88} gives
$Z_g = \sum_j S_{0j}^{2-2g}$ as the dimension of the space
of conformal blocks on a genus-$g$ curve with no insertions.
From the ordered chiral homology: the symmetric coinvariants
of the ordered chain complex compute the conformal blocks
via the factorization property (TUY~\cite{TUY89}). The
KZB local system at integrable level decomposes into $k+1$
channels labeled by integrable representations; each channel
$j$ contributes $S_{0j}^{2-2g}$ to the partition function.
The integrality of $Z_g$ follows from the identification
with the rank of the Verlinde bundle over~$\cM_g$.

\textup{(iv)} Handle attachment: a genus-$g$ surface with
a non-separating cycle cut open gives a genus-$(g{-}1)$
surface with two extra marked points. The bar complex sewing
map reglues these marked points, summing over the $k+1$
intermediate representations with the propagator
$S_{0j}^{-2}$ (the inverse of the $j$-th eigenvalue of
the sewing operator). This gives
$Z_{g+1} = \sum_j H_j \cdot S_{0j}^{2-2g}$.

\textup{(v)} Separating factorization: the bar complex
factorization under separating degeneration
(Theorem~A for the ordered bar complex on $\Ran^{\mathrm{ord}}$)
decomposes the genus-$g$ partition function into a sum over
channels of the product of the two lower-genus contributions,
normalized by the propagator $S_{0j}^{-2}$.

All numerical values in the table have been verified by
direct computation in
\texttt{verlinde\_ordered\_engine.py}
using three independent code paths
(direct $S$-matrix summation, quantum-dimension formula,
and handle recursion), cross-checked against the Verlinde
table of \texttt{verlinde\_shadow\_algebra.py}
($198$ tests, all passing).
\end{proof}

\begin{remark}[Verlinde bundle vs shadow tower]
\label{rem:verlinde-vs-shadow}
The Verlinde dimension $Z_g$ and the shadow partition
function $F_g = \kappa \cdot \lambda_g$ are different
invariants of the same chiral algebra. $Z_g$ is the rank of
the Verlinde bundle over $\cM_g$ (an integer, from the
$S$-matrix); $F_g$ is the first Chern class integrated
over $\overline{\cM}_g$ (a rational number, from $\kappa$).
For $V_k(\mathfrak{sl}_2)$:
$\kappa = 3(k{+}2)/4$ (equation~\eqref{eq:sl2-kappa}),
$F_1 = \kappa/24$, and $Z_1 = k + 1$.
The two invariants encode complementary data: $Z_g$ counts
the conformal blocks (combinatorial, from the modular tensor
category), while $F_g$ measures the curvature (analytic,
from the Hodge class $\omega_g = c_1(\lambda)$ on
$\overline{\cM}_g$; AP130: $\omega_g$ lives on
$\overline{\cM}_g$, not on the curve $\Sigma_g$).
\end{remark}

\begin{remark}[The ordered refinement]
\label{rem:verlinde-ordered-refinement}
The Verlinde dimension $Z_g = \dim(\int_{\Sigma_g}^{\mathrm{ord}}
V_k)_{\Sigma_n}$ is the dimension of the symmetric coinvariants.
The full ordered chiral homology
$\int_{\Sigma_g}^{\mathrm{ord}} V_k(\mathfrak{sl}_2)$
is strictly larger: its degree-$n$ component carries the
full KZB monodromy representation of the elliptic braid
group $\mathrm{EB}_n$
(Proposition~\ref{prop:ell-degree2},
Theorem~\ref{thm:ell-ordered-ch}), from which the Verlinde
dimension is recovered by taking $\Sigma_n$-coinvariants
and restricting to $n = 0$. The kernel $\ker(\av_n)$
carries the matrix-valued monodromy data: $R$-matrices,
$6j$-symbols, and (at genus~$\geq 1$) dynamical quantum
group data.
\end{remark}


% ================================================================

\section{Ordered chiral homology at genus~$2$:
$\int_{\Sigma_2}^{\mathrm{ord}} Y_\hbar(\mathfrak{sl}_2)$}
